\documentclass[12pt]{article}
\usepackage{sbc-template}
\usepackage[T1]{fontenc}
\usepackage[utf8]{inputenc}
\usepackage[brazil]{babel}
\usepackage{graphicx,url}
\sloppy
\title{Tempo de tela e indicadores de saúde mental em crianças e adolescentes: uma análise exploratória da NSCH de 2020 a 2024}
\author{Kauan Edson Leonel Lira\inst{1}, Eduarda Gonçalves Luiz de Andrade\inst{1}}
\address{Instituto Federal de São Paulo -- Câmpus Caraguatatuba -- Brasil}
\begin{document}
\maketitle
\begin{abstract}
This study examines parent-reported screen time and current anxiety and depression in US National Survey of Children's Health data. Five annual cross-sections yielded 166,343 records aged 6--17. Unweighted complete-case logistic models used 2020--2023 only. For four or more hours versus less than one, sample odds ratios were 1.79 for depression and 1.65 for anxiety. These exploratory estimates lack design-based uncertainty and cannot support population or causal inference.
\end{abstract}
\begin{resumo}
Este estudo examina tempo de tela e relatos de responsáveis sobre ansiedade e depressão atuais na National Survey of Children's Health dos EUA. Cinco edições transversais reuniram 166.343 registros de 6 a 17 anos. Regressões logísticas não ponderadas de casos completos usaram somente 2020--2023. Para quatro horas ou mais frente a menos de uma, as razões de chances amostrais foram 1,79 para depressão e 1,65 para ansiedade. As estimativas exploratórias não têm incerteza compatível com o desenho e não fundamentam inferência populacional ou causal.
\end{resumo}
\section{Introdução}
O tempo diante de dispositivos eletrônicos reúne práticas distintas, como assistir a programas, jogar e interagir em redes sociais. A possível relação dessas práticas com a saúde mental de crianças e adolescentes é relevante, mas uma medida agregada de duração não identifica conteúdo, contexto familiar ou motivos de uso. Na coorte norte-americana estudada por \cite{paulich2021}, associações com saúde mental e sono tiveram magnitudes modestas, enquanto relações sociais também apresentaram padrões distintos. A interpretação exige, portanto, distinguir duração, tipo de atividade e características dos participantes.

A \textit{National Survey of Children's Health} (NSCH) fornece dados anuais sobre crianças residentes em domicílios nos Estados Unidos \cite{censusdata}. A combinação de edições amplia o número de observações, mas não constitui seguimento longitudinal e exige tratamento do plano amostral \cite{censusguide}. Este trabalho pergunta como categorias de tempo de tela se relacionam com relatos de depressão e ansiedade atuais na amostra disponível. A hipótese exploratória é que categorias superiores apresentem maiores chances condicionais desses relatos, sem pressupor uma relação causal ou monotônica.

O objetivo é documentar a harmonização das edições de 2020 a 2024, examinar associações amostrais em 2020--2023 e relatar a avaliação exploratória temporal dos classificadores em 2024. A contribuição é um registro reproduzível das definições, dos resultados e das condições ainda necessárias para inferência populacional e validação preditiva. 
\section{Trabalhos relacionados}
A literatura associa alguns usos de tela a sintomas depressivos, mas os resultados variam segundo desenho e medida. Em acompanhamento de adolescentes canadenses, \cite{boers2019} distinguiram atividades digitais e associações entre pessoas e dentro da mesma pessoa. A metanálise de coortes prospectivas de \cite{gao2024} encontrou associação agregada entre exposição inicial e depressão posterior, acompanhada de heterogeneidade elevada; os estudos observacionais reunidos não eliminam confundimento.

Em contraste, \cite{li2025} encontraram associações longitudinais mais fracas que as transversais, sem associação longitudinal com ansiedade em sua análise de adolescentes australianos. As diferenças de população, instrumento e intervalo de acompanhamento impedem comparar diretamente coeficientes com a NSCH. A aparente divergência também pode refletir causalidade reversa ou fatores comuns ao uso e aos sintomas, hipóteses que dados transversais não conseguem separar.

A revisão longitudinal de \cite{tang2021} encontrou associações pequenas ou muito pequenas com sintomas depressivos posteriores e evidência limitada para ansiedade. Na revisão de revisões de \cite{valkenburg2022}, a maioria das sínteses qualificou as associações de redes sociais com saúde mental como fracas ou inconsistentes, embora algumas as interpretassem como substanciais. Em outra análise de adolescentes norte-americanos, \cite{kreski2021} não encontraram associação forte e consistente entre uso diário de redes sociais e sintomas depressivos. Essas medidas de redes sociais não equivalem ao tempo total de tela da NSCH.

Os instrumentos e o contexto influenciam a comparação dos achados. \cite{stiglic2019} identificaram evidência de qualidade variável em revisão de revisões sobre telas e saúde de crianças, enquanto \cite{orben2020} ressaltou que diferenças individuais e escolhas analíticas dificultam resumir experiências digitais em uma média única. Um estudo anterior com a própria NSCH examinou envolvimento com telas e funcionamento psicossocial \cite{przybylski2020}, mas usou edição e desfechos diferentes dos adotados aqui. A presente análise não estima exposição a plataformas específicas nem conteúdo consumido.

Sono ilustra a dificuldade de escolher covariáveis: revisão de \cite{hale2015} descreve associações frequentes entre telas e desfechos do sono, com limitações de mensuração e de causalidade. Se sono ou relações sociais forem intermediários, seu ajuste pode retirar parte de uma associação de interesse; se forem consequências do desfecho, a interpretação também muda. Assim, os modelos descritos abaixo representam associações condicionais a variáveis contemporâneas, não efeitos totais.
\section{Materiais e métodos}
\subsection{Fonte, população e harmonização}
Foram utilizados arquivos locais CSV do questionário \textit{topical} das edições de 2020 a 2024 e o dataset Parquet produzido pelo projeto. Cada edição é transversal e a unidade é a criança selecionada; registros agrupados não identificam trajetórias individuais. A análise harmonizada restringiu idade a 6--17 anos, universo comum às variáveis usadas. A rotina verificou identificador, idade, formulário, estrato, peso e domínios; comparou frequências de tela e totais e pesos por sexo com publicações oficiais. Esses controles sustentam compatibilidade, mas não autenticam os arquivos originais byte a byte.

Os arquivos de 2020 e 2021 têm nome da versão aprimorada e seus controles de pesos coincidiram com frequências publicadas conforme o relatório do projeto. O guia oficial exige versões aprimoradas das edições anteriores a 2022 em comparações com anos posteriores \cite{censusguide}; o nome, isoladamente, não comprova a procedência. As listas de variáveis e as regras por ano estão documentadas em \texttt{docs/harmonizacao.md}. A comparação de pesos publicados não substitui conferência independente de proveniência completa.

\subsection{Exposição, desfechos e dados ausentes}
A variável \texttt{SCREENTIME} descreve tempo habitual em dias de semana em cinco categorias: menos de uma hora, uma, duas, três e quatro horas ou mais. A última é aberta e não representa duração exata. A lista de 2024 exclui explicitamente tarefas escolares \cite{census2024}; segundo a comparação documental do projeto, essa frase não aparece no item das listas anteriores, o que limita a suposição de equivalência perfeita com 2024 \cite{census2020}. A medida não separa dispositivos, conteúdo ou uso simultâneo.

Depressão e ansiedade foram construídas separadamente com os itens de condição alguma vez informada por profissional e condição atual (\texttt{K2Q32A/B} e \texttt{K2Q33A/B}). Histórico negativo com salto documentado implica condição atual negativa; histórico positivo exige resposta válida ao item atual. Histórico desconhecido, atual indeterminado ou combinações inconsistentes permanecem ausentes. Trata-se de relato do responsável sobre condição informada, não de avaliação clínica conduzida neste estudo. Na regressão foram removidos registros com ausência em qualquer variável do respectivo modelo; não houve imputação.

\subsection{Análise e desenho amostral}
O código existente ajusta regressões logísticas separadas, sem penalização, com intercepto, idade contínua e indicadores para sexo, sono, atividade física, dificuldade para fazer amigos, ano e tempo de tela. A referência é menos de uma hora. As razões de chances (\textit{odds ratios}, OR) são exponenciais dos coeficientes condicionais; não são razões de prevalência. O algoritmo \texttt{lbfgs} usa tolerância de $10^{-8}$ e limite de 3.000 iterações, com verificação de convergência. As regressões de associação usam exclusivamente 2020--2023. Não foram executados testes de tendência ou comparações por idade.

As regressões não usam \texttt{FWC}, estratos nem erro-padrão adequado ao desenho. Portanto, seus coeficientes descrevem apenas os casos completos da amostra e não são estimativas nacionais; tampouco há intervalos de confiança ou testes de hipótese. Para estimativa populacional futura, o guia do Census recomenda pesos revisados e estimação de variância compatível com o plano; para tamanho populacional anual médio de cinco anos, divide-se o peso anual por cinco, o que não altera proporções por si só \cite{censusguide}. Seus exemplos usam estado e estrato, unidade domiciliar e procedimentos próprios para pesquisas complexas. A definição de estratos e identificadores entre anos e a estimação por domínio devem ser validadas antes de produzir intervalos: simplesmente ponderar uma regressão convencional não resolve a variância.

Os dados são arquivos públicos disponibilizados pelo Census; o processamento usa identificadores apenas para integridade e não apresenta registros individuais. A aprovação ética aplicável à instituição e ao evento deve ser verificada pelos responsáveis. O pipeline e versões estão registrados no material de reprodução.

\subsection{Experimentos preditivos}
Em análise complementar, o projeto comparou dois conjuntos de preditores: A, com tempo de tela, e B, sem ele; ambos incluem idade, sexo, sono, atividade física e dificuldade para fazer amigos. Em 2020--2023, regressão logística, floresta aleatória e XGBoost foram comparados a um classificador de referência por validação cruzada estratificada em cinco partes, com seleção por área sob a curva ROC (ROC AUC). A transformação categórica e o tratamento de faltantes integram os pipelines; para variáveis numéricas usa-se imputação pela mediana e para categóricas, categoria de ausência. O limiar maximiza F1 em previsões fora do respectivo fold de desenvolvimento. Como parâmetros foram escolhidos usando esses mesmos folds, tais previsões não constituem validação aninhada independente.

Os modelos salvos foram avaliados em registros de 2024 com desfecho conhecido. O projeto também calculou intervalos percentílicos por 1.000 reamostras bootstrap pareadas de registros para a diferença A menos B; esses intervalos não consideram pesos nem estratos e condicionam-se aos modelos ajustados. Após uma primeira avaliação de 2024, o protocolo de treinamento e limiar foi revisado; portanto, 2024 deixou de ser uma avaliação totalmente cega. Os escores não foram calibrados e a análise mede desempenho na amostra, sem equivaler a diagnóstico ou efeito causal da variável incluída.
\section{Resultados}
A Tabela~\ref{tab:fluxo} apresenta as contagens por edição antes e depois da seleção etária. O conjunto final reuniu 166.343 registros, sem duplicação do identificador analítico. Havia 2.705 valores ausentes de tempo de tela, 965 de depressão atual e 1.642 de ansiedade atual; essas ausências podem se sobrepor. Em todo o conjunto etário, os rótulos de depressão compreenderam 11.193 positivos e 154.185 negativos; os de ansiedade, 26.051 positivos e 138.650 negativos. Essas contagens são amostrais, sem ponderação.
\begin{table}[ht]
\centering
\caption{Registros por edição, antes e após seleção etária. Contagens não ponderadas.}
\label{tab:fluxo}
\begin{tabular}{lrr}\hline
Ano & Arquivo original & 6--17 anos \\\hline
2020 & 42.777 & 30.636 \\
2021 & 50.892 & 30.173 \\
2022 & 54.103 & 34.362 \\
2023 & 55.162 & 33.638 \\
2024 & 51.375 & 37.534 \\\hline
\end{tabular}
\end{table}

Nas quatro edições de desenvolvimento havia 128.809 registros de 6--17 anos. Após exclusão por dados ausentes no desfecho ou em variáveis do modelo, restaram 124.795 para depressão e 124.299 para ansiedade, com exclusão respectiva de 4.014 e 4.510 registros. A Tabela~\ref{tab:or} mostra as OR condicionais obtidas na reprodução. Para quatro horas ou mais, as OR foram 1,79 e 1,65, respectivamente; na categoria de uma hora, a OR de depressão foi 0,97. Não há evidência de precisão estatística nessas saídas, nem justificativa para interpretar diferenças pequenas entre categorias como conclusivas.
\begin{table}[ht]
\centering
\caption{Razões de chances condicionais amostrais, 2020--2023, por tempo de tela.}
\label{tab:or}
\begin{tabular}{lrr}\hline
Tempo de tela & Depressão & Ansiedade \\\hline
Menos de 1 hora & 1,00 & 1,00 \\
1 hora & 0,97 & 1,13 \\
2 horas & 1,22 & 1,27 \\
3 horas & 1,37 & 1,43 \\
4 horas ou mais & 1,79 & 1,65 \\\hline
\end{tabular}

\smallskip\footnotesize Referência: menos de uma hora. Casos completos; ajuste por idade, sexo, sono, atividade física, amizade e ano. Sem pesos ou intervalos de confiança.
\end{table}

A Figura~\ref{fig:or} apresenta os mesmos coeficientes para facilitar a comparação visual. Os pontos não incluem barras de incerteza porque os modelos de associação não as estimaram; a posição de um ponto não demonstra diferença estatisticamente estabelecida.
\begin{figure}[ht]
\centering
\includegraphics[width=0.95\textwidth]{figuras/associacao.png}
\caption{Razões de chances condicionais não ponderadas em 2020--2023; sem intervalos de confiança.}
\label{fig:or}
\end{figure}

\subsection{Avaliação temporal exploratória}
Para cada desfecho e conjunto de preditores, XGBoost foi selecionado pela ROC AUC no desenvolvimento. A Tabela~\ref{tab:predicao} resume a avaliação em 2024. A área ROC mede a ordenação dos escores; a precisão média (AP) resume a curva precisão--sensibilidade e depende da frequência do desfecho. Precisão é a fração de positivos previstos que corresponde a positivos observados; sensibilidade é a fração dos positivos observados identificada pelo limiar. Ambas e F1 dependem do limiar escolhido.
\begin{table}[ht]
\centering
\caption{Avaliação exploratória de XGBoost em 2024, sem ponderação.}
\label{tab:predicao}
\begin{tabular}{llrrrrr}\hline
Desfecho & Modelo & ROC AUC & AP & Precisão & Sensib. & F1 \\\hline
Depressão & A & 0,839 & 0,327 & 0,300 & 0,508 & 0,377 \\
Depressão & B & 0,835 & 0,322 & 0,320 & 0,471 & 0,381 \\
Ansiedade & A & 0,761 & 0,450 & 0,402 & 0,546 & 0,463 \\
Ansiedade & B & 0,758 & 0,448 & 0,403 & 0,541 & 0,462 \\\hline
\end{tabular}

\smallskip\footnotesize A inclui tempo de tela; B o exclui. Limiar escolhido no desenvolvimento e diferente entre modelos. Desfecho conhecido: depressão $n=37.275$; ansiedade $n=37.128$. AP: precisão média.
\end{table}

Em 2024, os registros positivos somaram 2.537 para depressão e 6.450 para ansiedade. Na comparação pareada A menos B, as diferenças de ROC AUC foram 0,00425 para depressão e 0,00291 para ansiedade; os intervalos percentílicos bootstrap de 95\% foram, respectivamente, [0,00268; 0,00571] e [0,00140; 0,00438]. Esses intervalos descrevem a variabilidade da reamostragem de registros sob modelos fixos e não representam incerteza do plano amostral, da seleção de modelos ou da adaptação do protocolo. As diferenças são pequenas, e F1 diminuiu para depressão e aumentou pouco para ansiedade. Não se deve interpretar esse contraste preditivo como efeito isolado de telas sobre a saúde.

\section{Discussão}
As OR amostrais das categorias mais altas foram superiores à referência nos dois modelos, mas não constituem uma estimativa de associação para a população dos Estados Unidos. A diferença entre os dois desfechos e a proximidade da unidade para depressão na categoria de uma hora desaconselham resumir o resultado como uma progressão comprovada. Sem intervalos de confiança que incorporem o desenho, não é possível avaliar a precisão nem concluir que categorias diferem entre si.

O padrão exploratório é compatível com a existência de associações relatadas em estudos anteriores, mas os desfechos e as populações diferem daqueles de \cite{boers2019,paulich2021}. A atenuação longitudinal descrita por \cite{li2025} reforça que associação transversal não determina direção temporal. A metanálise de \cite{gao2024} também não elimina heterogeneidade nem fornece um coeficiente diretamente transferível para a escala categórica da NSCH. O relato do responsável pode depender de reconhecimento de sintomas, acesso a serviços e recordação, enquanto a exposição não mede a qualidade da experiência digital.

Além da seleção por não resposta, os casos completos podem diferir dos excluídos. Idade e sexo não esgotam fatores socioeconômicos, familiares ou de saúde preexistente. Sono, atividade e amizade podem ser confundidores, mediadores ou consequências e foram medidos contemporaneamente; ajustar por eles não identifica uma via causal. A faixa de 6--17 anos mistura etapas de desenvolvimento, e a categoria aberta de exposição impede distinguir durações mais longas. Diferenças na redação entre edições e mudanças no contexto do período também restringem comparações temporais. A regressão de associação não inclui 2024 e não descreve a associação nesse ano.

No componente preditivo, uma ROC AUC relativamente alta pode coexistir com precisão limitada quando há desequilíbrio entre classes. No modelo A para depressão, por exemplo, aproximadamente metade dos positivos observados foi identificada, mas menos de um terço das classificações positivas correspondia ao relato positivo na amostra. O desempenho é específico ao conjunto de variáveis e ao relato dos responsáveis; a dificuldade de fazer amigos pode conter informação contemporânea relacionada ao próprio desfecho. O resultado não estabelece utilidade clínica, calibração ou transporte para crianças brasileiras. Além disso, a revisão do protocolo após contato prévio com 2024 requer uma avaliação externa ou temporal inteiramente nova para validação independente.

Uma extensão defensável exige verificar a procedência enhanced, documentar comparabilidade dos instrumentos e especificar o desenho completo antes da seleção de domínio. Deverá apresentar contagens não ponderadas, prevalências ponderadas, intervalos de confiança apropriados e análises de sensibilidade para ausências, idade e conjuntos de ajuste justificados. A escolha de covariáveis deve refletir uma pergunta causal explícita caso se pretenda estimar efeitos, e não ser ditada apenas pela disponibilidade de campos.
\section{Conclusão}
Na amostra de 2020--2023, os modelos exploratórios encontraram maiores chances condicionais de relatos de depressão e ansiedade nas categorias superiores de tempo de tela em comparação com menos de uma hora. As estimativas são não ponderadas e sem medidas adequadas de incerteza; não respondem se a associação ocorre na população nem se o uso de tela causa condições de saúde mental. A harmonização de 2020--2024 e a reprodução dos coeficientes fornecem base para uma análise futura que respeite o desenho amostral, examine a comparabilidade e relate a incerteza. Nenhum resultado constitui diagnóstico médico.
\bibliographystyle{sbc}
\bibliography{referencias}
\end{document}
